\chapter{Линейная регрессия, диагностика и обобщённые линейные модели}
\label{ch:regression}

Главы~\ref{ch:basics}--\ref{ch:multiple} заложили аппарат распределений, доверительных интервалов и проверки гипотез. Мы научились спрашивать, отличается ли среднее двух групп и согласуются ли данные с $H_0$, опираясь на $t$-, $F$- и $\chi^2$-критерии при известных предпосылках.
Регрессионный анализ переносит эту логику на задачу описания зависимости. Как изменяется отклик $y$ при варьировании предикторов $x_1,\ldots,x_k$.
Вместо сравнения двух выборочных средних мы оцениваем параметры линейной (или обобщённой линейной) модели и проверяем, какие из них статистически значимы при заданном уровне $\alpha$.

Метод наименьших квадратов (МНК) даёт линейную аппроксимацию условного матожидания $\mathbb{E}(y\mid x)$ \citep{casella2002}. Обобщённые линейные модели (GLM) расширяют этот аппарат на бинарные и счётные отклики, где нормальность остатков уже неуместна, а связь между $\mathbb{E}(y\mid x)$ и $x^\top\beta$ задаётся через связующую функцию.
Без проверки предположений и содержательной интерпретации коэффициентов регрессия легко превращается в формальную подгонку чисел. $R^2$ растёт, p-value выглядят убедительно, а прогноз за пределами обучающей выборки оказывается бессмысленным.
Структура главы повторяет логику предыдущих. Сначала модель и оценивание, затем инференс (доверие к p-value при выполнении предпосылок), диагностика (проверка этих предпосылок) и обобщение на другие типы отклика.
Настоящая глава связывает оценивание, инференс и диагностику в единую цепочку. От постановки МНК через классические $t$- и $F$-тесты к проверке остатков и, наконец, к GLM для качественных и счётных данных.

\section{Постановка и метод наименьших квадратов}

Пусть $n$ объектов описываются матрицей предикторов $X\in\mathbb{R}^{n\times(k+1)}$ (первый столбец~--- единицы) и вектором отклика $y\in\mathbb{R}^n$.
Строка $x_i^\top$~--- признаки $i$-го объекта.
Интуитивно каждый $y_i$~--- реализация случайной величины с условным распределением, зависящим от $x_i$. Регрессия ищет простое правило, которое наилучшим образом предсказывает типичное значение $y_i$ при заданном $x_i$.

\begin{definition}
Линейная модель задаёт отклик как сумму линейной комбинации предикторов и случайной ошибки:
\[ y = X\beta + \varepsilon, \qquad \mathbb{E}(\varepsilon\mid X)=0.
\]
\end{definition}

\begin{definition}
Метод наименьших квадратов (МНК, OLS) минимизирует $\|y-X\beta\|_2^2$ и эквивалентен минимизации $\sum_i (y_i-\mathbb{E}(y_i\mid x_i))^2$ в классе линейных предикторов. Оценка $\hat\beta$ называется OLS-оценкой.
\end{definition}

Геометрически $\hat y=X\hat\beta$~--- ортогональная проекция $y$ на столбцовое пространство $X$. Остатки $y-\hat y$ перпендикулярны всем столбцам $X$, что и выражает условие первого порядка.
Условие первого порядка $2X^\top(y-X\beta)=0$ даёт нормальные уравнения $X^\top X\beta=X^\top y$ и при $\rank X=k+1$:
\[ \hat\beta = (X^\top X)^{-1}X^\top y, \qquad \hat y = Hy, \quad H=X(X^\top X)^{-1}X^\top.
\]
Матрица $H$ (hat matrix) проецирует $y$ на столбцовое пространство $X$. Диагональные элементы $h_{ii}$ измеряют рычаг $i$-го наблюдения~--- насколько сильно точка $(x_i,y_i)$ тянет подгонку к себе.
Следствие $H^2=H$. Повторная проекция не меняет результат. Сумма диагональных элементов $\sum_i h_{ii}=k+1$ связывает общий рычаг с числом параметров модели.

Разложение сумм квадратов $TSS=\sum(y_i-\bar y)^2$, $ESS=\sum(\hat y_i-\bar y)^2$, $RSS=\sum(y_i-\hat y_i)^2$, $TSS=ESS+RSS$.
Коэффициент детерминации $R^2=ESS/TSS=1-RSS/TSS$.
Приведённый $R^2_a$ $1-(1-R^2)\frac{n-1}{n-k-1}$~--- для сравнения моделей с разным числом регрессоров.
Следует помнить. Высокий $R^2$ не доказывает причинность и не гарантирует хороший прогноз на новых данных. Он лишь количественно описывает долю вариации $y$, объяснённую линейной комбинацией $x$ на данной выборке \citep{casella2002}.

\begin{definition}
(1)~линейность $y=X\beta+\varepsilon$. (2)~случайная выборка. $(x_i,y_i)$ независимы. (3)~полный ранг $\rank X=k+1$. (4)~$\mathbb{E}(\varepsilon\mid X)=0$. (5)~гомоскедастичность $\mathbb{D}(\varepsilon\mid X)=\sigma^2 I$.
При их выполнении МНК-оценки $\hat\beta$ несмещены, состоятельны и имеют наименьшую дисперсию среди линейных по $y$ несмещённых оценок (теорема Гаусса--Маркова).
\end{definition}

Предположения (2) и (4) связывают регрессию с i.i.d.-посылкой глав~1--2. Нарушение независимости (например, временные ряды или кластеры) требует робастных или специализированных процедур, которые мы затронем кратко и подробно разберём в гл.~\ref{ch:timeseries}.

Эскиз доказательства теоремы Гаусса--Маркова. Пусть $\tilde\beta=Ay$~--- произвольная линейная несмещённая оценка $\beta$ ($A$ не зависит от $y$, $AX=I_{k+1}$).
Тогда $\tilde\beta-\hat\beta=(A-(X^\top X)^{-1}X^\top)y$ не зависит от истинного $\beta$, а при $\mathbb{E}(\varepsilon\mid X)=0$ и $\mathbb{D}(\varepsilon\mid X)=\sigma^2 I$:
\[ \cov(\tilde\beta-\hat\beta\mid X) =\sigma^2\bigl(A-(X^\top X)^{-1}X^\top\bigr) \bigl(A-(X^\top X)^{-1}X^\top\bigr)^\top.
\]
Матрица $\cov(\hat\beta\mid X)=\sigma^2(X^\top X)^{-1}$.
Разность ковариаций $\cov(\tilde\beta\mid X)-\cov(\hat\beta\mid X)$ положительно полуопределена. Это следует из тождества $(A-(X^\top X)^{-1}X^\top)X=0$ и разложения $\cov(\tilde\beta\mid X)=\cov(\tilde\beta-\hat\beta\mid X)+\cov(\hat\beta\mid X)$.
Значит, $\hat\beta$ оптимален в классе линейных несмещённых оценок.
Замечание. Нелинейные или смещённые оценки (ridge, lasso) могут дать меньшую среднеквадратичную ошибку при мультиколлинеарности.
В терминах гл.~\ref{ch:basics} теорема отвечает на вопрос, какую линейную оценку выбрать, но не на вопрос, верна ли линейная модель вообще. Для этого нужна диагностика остатков.
Теорема не утверждает, что МНК~--- лучший метод вообще. Она ограничивает класс сравнения линейными несмещёнными оценками, что согласуется с частотной логикой глав~1--2, где мы ищем процедуры с контролируемым уровнем ошибки I рода при корректных предпосылках.

Дисперсия $\hat\beta_j$ растёт с $\sigma^2$ и мультиколлинеарностью $x_j$:
\[ \mathbb{D}(\hat\beta_j\mid X)=\frac{\sigma^2}{\mathrm{TSS}_j(1-R_j^2)}, \]
где $R_j^2$~--- $R^2$ регрессии $x_j$ на остальные столбцы $X$, $\mathrm{TSS}_j=\sum_i(x_{ij}-\bar x_j)^2$.
В матричной форме $\mathbb{D}(\hat\beta\mid X)=\sigma^2(X^\top X)^{-1}$. Почти линейно зависимые столбцы делают $X^\top X$ плохо обусловленной, и малые изменения $y$ могут сильно сдвигать $\hat\beta$. Это не ошибка метода, а сигнал о недостатке информации в данных.

VIF (variance inflation factor) для $j$-го предиктора:
\[ \mathrm{VIF}_j=\frac{1}{1-R_j^2}.
\]
$\mathrm{VIF}_j>10$ (иногда $>5$) сигнализирует о сильной мультиколлинеарности. Доверительные интервалы для $\beta_j$ расширяются примерно в $\sqrt{\mathrm{VIF}_j}$ раз. Лечение~--- отбор признаков, центрирование, регуляризация или содержательное объединение коррелированных факторов.
Важно не путать статистическую и содержательную проблему. Высокий VIF не доказывает, что один из предикторов лишний с точки зрения причинности, но предупреждает, что отдельные $t$-тесты для коэффициентов могут быть нестабильны при малом сдвиге данных.

Пример. Внешность и оценка курса. В исследовании красота и оценка курса измерены внешность преподавателя (\texttt{looks}) и итоговую оценку курса (\texttt{course\_eval}).
Простая линейная регрессия показывает положительную связь. На каждую единицу прироста \texttt{looks} ожидается изменение $\hat\beta_1$ баллов оценки при прочих равных (здесь других предикторов нет).
Интерпретация при прочих равных имеет смысл только в рамках линейной модели и при отсутствии смешивающих факторов, не включённых в $X$.
Даже при значимом $\hat\beta_1$ нельзя автоматически заключать причинность. Внешность может коррелировать с другими характеристиками преподавателя, не измеренными в данных.
График остатков (рис.~\ref{fig:reg-residuals}) позволяет проверить линейность и гомоскедастичность~--- первый шаг диагностики, без которого $t$- и $F$-критерии следующего раздела теряют обоснованность.

\begin{figure}[htbp]
    \centering
    \includegraphics[width=\FigWidth]{\FigDir/chapter-5/goodres.png}
    \caption{OLS-подгонка и остатки (иллюстрация на синтетических данных).}
    \label{fig:reg-residuals}
\end{figure}

\section{Инференс в линейной модели}

При нормальности ошибок $\varepsilon\sim N(0,\sigma^2 I)$ (шестое предположение) МНК совпадает с МП. $\hat\beta\sim N(\beta,\sigma^2(X^\top X)^{-1})$, $\hat\sigma^2=RSS/(n-k-1)$, $(n-k-1)\hat\sigma^2/\sigma^2\sim\chi^2_{n-k-1}$.
Шестое предположение не входит в теорему Гаусса--Маркова, но необходимо для точных (не только асимптотических) интервалов и p-value в конечных выборках~--- ровно так же, как $t$-критерий в гл.~\ref{ch:basics} опирается на нормальность при малых $n$.
При больших $n$ центральная предельная теорема оправдывает асимптотические интервалы даже при умеренном отклонении от нормальности ошибок, но на малых выборках и при тяжёлых хвостах $y$ Q--Q график остатков остаётся обязательным.
Оценка $\hat\sigma^2=RSS/(n-k-1)$ несмещена при гомоскедастичности. Подстановка в $t$-статистику даёт распределение Стьюдента с $n-k-1$ степенями свободы~--- прямой аналог одновыборочного $t$-критерия, но для линейной комбинации коэффициентов $\hat\beta$.
Для $c\in\mathbb{R}^{k+1}$:
\[ \frac{c^\top(\hat\beta-\beta)}{\hat\sigma\sqrt{c^\top(X^\top X)^{-1}c}} \sim t_{n-k-1}.
\]
Отсюда $t$-интервалы для $\beta_j$ и для $\mathbb{E}(y_0\mid x_0)=x_0^\top\beta$:
\[ \hat\beta_j \pm t_{1-\alpha/2}\,\hat\sigma\sqrt{(X^\top X)^{-1}_{jj}}, \qquad x_0^\top\hat\beta \pm t_{1-\alpha/2}\,\hat\sigma\sqrt{x_0^\top(X^\top X)^{-1}x_0}.
\]
Предсказательный интервал для нового $y_0=x_0^\top\beta+\varepsilon_0$ шире на $\hat\sigma$ (под корнем появляется $+1$):
\[ x_0^\top\hat\beta \pm t_{1-\alpha/2}\,\hat\sigma\sqrt{1+x_0^\top(X^\top X)^{-1}x_0}.
\]
Различие доверительного и предсказательного интервалов принципиально. Первый отвечает на вопрос о среднем отклике при $x_0$, второй~--- о конкретном будущем наблюдении с учётом шума $\varepsilon_0$.
На практике путаница между ними приводит к завышенной уверенности в точечных прогнозах.
Если $x_0$ лежит далеко от центра облака предикторов (большой $x_0^\top(X^\top X)^{-1}x_0$), оба интервала резко расширяются~--- экстраполяция статистически и содержательно рискованна.
Различие интервала для $\mathbb{E}(y\mid x_0)$ и интервала для нового $y_0$ часто путают в отчётах. Первый отвечает, где проходит линия регрессии в среднем. Второй~--- куда попадёт следующее наблюдение.
На практике предсказательный интервал шире. Использование доверительного интервала для прогноза завышает уверенность в точечных прогнозах.

Частичный $F$-критерий. Разобьём $X=(X_1,X_2)$, $\beta^\top=(\beta_1^\top,\beta_2^\top)$, где $X_2$ содержит $k_1$ столбцов, подлежащих совместной проверке.
Частичный (nested) $F$-критерий для $H_0\colon\beta_2=0$:
\[ F=\frac{(RSS_r-RSS_{ur})/k_1}{RSS_{ur}/(n-k-1)} \sim F_{k_1,n-k-1}\quad\text{при }H_0, \]
$RSS_r$~--- остаточная сумма квадратов редуцированной модели (только $X_1$), $RSS_{ur}$~--- полной.
При $k_1=1$ $F$ эквивалентен двустороннему $t$-критерию для одного $\beta_j$.
При $k_1=k$ получаем общий $F$ для $H_0\colon\beta_1=\cdots=\beta_k=0$. $F=\frac{R^2/k}{(1-R^2)/(n-k-1)}$.
Несогласованность частичного $F$ и отдельных $t$-тестов часто объясняется мультиколлинеарностью или множественными сравнениями~--- напоминание о процедурах контроля FWER/FDR из гл.~\ref{ch:multiple}.
Совместная проверка группы регрессоров предпочтительна, когда они содержательно образуют единый блок (например, набор dummy для фактора).
Общий $F$ для $H_0\colon\beta_1=\cdots=\beta_k=0$ отвечает на вопрос, есть ли вообще линейная связь $y$ с $X$. Отдельные $t$-тесты спрашивают, значим ли каждый коэффициент при прочих фиксированных.

Категориальные предикторы кодируют dummy-переменными (базовый уровень~--- референс). One-hot с константой даёт мультиколлинеарность.
Категориальный фактор с $m$ уровнями включают или исключают целиком ($m-1$ dummy).
Пошаговая регрессия (вперёд/назад) удобна для отбора, но множественные сравнения требуют осторожности. Эксперимент Фридмана (1983) показывает риск ложных включений при $n\approx k$. На чистом шуме пошаговый отбор находит значимые предикторы, а $t$- и $F$-критерии, посчитанные на той же выборке, становятся антиконсервативны.
Сравнение невложенных моделей.~--- критерий Дэвидсона--Маккиннона $J$.

Критерий Дэвидсона--Маккиннона. Пусть конкурируют модели $M_1$ и $M_2$ с предсказаниями $\hat y$ и $\hat{\hat y}$.
Вложите прогноз соперника как регрессор. $M_1$~--- $y=X_1\beta_1+\gamma_1\hat{\hat y}+\varepsilon$, $M_2$~--- $y=X_2\beta_2+\gamma_2\hat y+\varepsilon$.
Проверьте $H_0\colon\gamma_1=0$ и $H_0\colon\gamma_2=0$ (критерий Стьюдента или $F$).
Если обе $H_0$ принимаются, модели неразличимы. Если отвергается только $\gamma_1=0$, $M_1$ предпочтительнее, и наоборот. Если обе отвергнуты, обе спецификации неполны.
Критерий $J$ переносит идею вложенного $F$-теста на случай, когда модели не вкладываются друг в друга~--- типичная ситуация при выборе между линейной и логарифмической спецификацией.
Практически сначала проверяют вложенные гипотезы (частичный $F$, критерий $G$ в GLM), и только при необходимости~--- $J$ для конкурирующих спецификаций.
Модель с одним предиктором даёт значимый $t$-тест для $\beta_1$. Добавление коррелированного $x_2$ раздувает SE($\hat\beta_1$) и ослабляет значимость~--- не потому что эффект исчез, а потому что мультиколлинеарность размывает вклад. Частичный $F$ на блок $(x_1,x_2)$ отвечает, нужны ли оба вместе. Это содержательнее двух отдельных $t$-тестов при коррелированных предикторах (см.~гл.~\ref{ch:multiple}).

\section{Диагностика регрессии}

Остатки $\hat\varepsilon_i=y_i-\hat y_i$. Стандартизованные $\tilde\varepsilon_i=\hat\varepsilon_i/\hat\sigma$.
Строят графики $\tilde\varepsilon$ от $\hat y$, от каждого $x_j$ и от индекса $i$.
Графический анализ дополняет формальные тесты. Глаз часто замечает систематические паттерны (воронка, кривизна), которые на малых $n$ не достигают значимости в критериях.
Диагностика~--- не формальность перед публикацией, а проверка того, что частотные гарантии $t$- и $F$-критериев (гл.~\ref{ch:basics}) применимы к конкретным данным.
Нарушение гомоскедастичности не делает $\hat\beta$ смещённым, но искажает стандартные ошибки и, следовательно, p-value и интервалы.

Линейность. Систематическая кривизна на графике остатков~--- добавить нелинейные члены или преобразование Бокса--Кокса для $y$.
Гомоскедастичность. Критерий Бройша--Пагана ($LM=nR^2$ регрессии $\hat\varepsilon_i^2$ на $X$, $\chi^2_k$). При нарушении~--- робастные стандартные ошибки Уайта (HC0--HC4) или взвешенный МНК.
Нормальность. Q--Q график, критерий Шапиро--Уилка (на малых $n$).
Влиятельные точки. Рычаг $h_{ii}$, расстояние Кука $D_i$. Большие $D_i$ и большие $|\tilde\varepsilon_i|$ требуют содержательной проверки.

Рычаг и расстояние Кука. Диагональ $h_{ii}$ матрицы $H$ лежит в $(0,1)$, $\sum_i h_{ii}=k+1$.
Большие $h_{ii}$ означают, что наблюдение далеко от центра облака $X$ (потенциально высокий рычаг).
Расстояние Кука измеряет суммарное влияние $i$-го пункта на все $\hat y_j$:
\[ D_i=\frac{\sum_{j=1}^n(\hat y_j-\hat y_{j(i)})^2}{(k+1)\hat\sigma^2} =\frac{\hat\varepsilon_i^2}{(k+1)\hat\sigma^2}\cdot\frac{h_{ii}}{(1-h_{ii})^2}, \]
где $\hat y_{j(i)}$~--- предсказание модели, обученной без $i$-го наблюдения.
Пороги. $D_i>1$, $D_i>4/n$, $D_i>3\bar D$ или визуальный анализ. Сочетайте с анализом остатков, а не удаляйте точки автоматически.
Точка может быть влиятельной без быть выбросом по $y$ (большой $h_{ii}$, малый остаток) и наоборот~--- содержательный контекст решает, является ли наблюдение ошибкой измерения или закономерным экстремумом.
Диагностика~--- не удаление выбросов до регрессии, а понимание механизма. Большое $D_i$ при малом остатке означает, что точка тянет наклон, но не противоречит модели. Большой остаток при умеренном $h_{ii}$~--- локальная аномалия $y$. Полный отчёт перечисляет оба типа точек и их влияние на $\hat\beta$, а не только $R^2$ до и после удаления.

Преобразование Бокса--Кокса. Для положительного $y$ семейство преобразований стабилизирует дисперсию:
\[ y^{(\lambda)}=
\begin{cases}
\dfrac{y^\lambda-1}{\lambda}, & \lambda\neq 0,\\[6pt] \ln y, & \lambda=0.
\end{cases}
\]
Подбор $\lambda$ по сетке на $[-2,2]$. Для каждого $\lambda$ строят регрессию на $y^{(\lambda)}$ и минимизируют приведённый $RSS(\lambda)$ (с поправкой Якобиана $\prod_i y_i^{\lambda-1}$ для $W$).
Доверительный интервал для $\lambda$~--- значения, где кривая $RSS(\lambda)$ не превышает $\min RSS\cdot \exp(\chi^2_{1,1-\alpha}/n)$. Если $1\in$ интервал, преобразование может не понадобиться.
Практически выбирают удобное $\lambda$ (0, $\pm\frac12$, 1).
Интерпретация коэффициентов после преобразования меняется. На шкале $\log y$ прирост предиктора интерпретируют как относительное изменение отклика, что часто содержательнее для положительных величин (доход, цена, объём).
Связь с гл.~\ref{ch:timeseries}. То же семейство применяют для стабилизации дисперсии временных рядов перед моделированием ARIMA.

\begin{itemize}
    \item Исключение неудобных точек без документирования смещает выводы.
    \item $R^2$ не измеряет качество прогноза на новых данных и растёт
          при любом добавлении регрессора.
    \item Корреляция предикторов не означает, что один из них лишний с точки
          зрения причинности, но доверительные интервалы расширяются.
    \item Экстраполяция за пределы диапазона $X$ опасна. Модель не проверена
          в новой области.
    \item Пошаговый отбор и классические p-value на той же выборке
          завышают значимость (Фридман).
\end{itemize}

Пример. Число визитов к врачу. В данных о визитах к врачу учтено число визитов (\texttt{DVISITS}) и демографические/медицинские предикторы.
Poisson-подобный отклик с избытком нулей мотивирует переход к GLM (раздел~\ref{sec:glm-ch5}) или моделям с hurdle. Линейная регрессия для счётного $y$ даёт предсказания вне $[0,\infty)$ и игнорирует дискретность распределения.
Типичная картина~--- положительный наклон на графике остатков от $\hat y$, завышенная дисперсия при больших $\hat\mu$~--- признаки неадекватности OLS.

Пример. Переломы и риск (GLOW500). В исследовании GLOW (переломы) бинарный отклик \texttt{FRACTURE} и факторы риска (\texttt{PRIORFRAC}, \texttt{BMI}, \texttt{SMOKE} и др.).
Логистическая регрессия оценивает $\logit P(y=1\mid x)=\log\frac{\pi}{1-\pi}=x^\top\beta$.
Бинарный отклик не может быть смоделирован нормальными остатками OLS. GLM задаёт правильное распределение и связь между вероятностью и предикторами.

\section{Обобщённые линейные модели}
\label{sec:glm-ch5}

\begin{definition}
GLM задаёт $\mathbb{E}(y\mid X)=\mu$ через связующую функцию $g$:
\[ g(\mu)=X\beta, \qquad \mu=g^{-1}(X\beta).
\]
Распределение $y$~--- из экспоненциального семейства:
\[ f(y,\theta,\phi)=\exp\!\left(\frac{y\theta-b(\theta)}{a(\phi)}+c(y,\phi)\right).
\]
Оценка $\hat\beta$~--- численный МП (Ньютон--Рафсон, scoring Фишера). Асимптотически $\hat\beta\sim N(\beta,I^{-1}(\beta))$, где $I(\beta)$~--- информационная матрица Фишера.
\end{definition}

OLS~--- частный случай GLM с нормальным распределением и тождественной связью. Логистическая и Poisson-регрессия~--- наиболее частые прикладные расширения.
Идея линейного предиктора $X\beta$ сохраняется. Меняются распределение $y$ и связь между $\mu=\mathbb{E}(y\mid x)$ и $X\beta$.
Таким образом, GLM~--- не отдельная дисциплина, а обобщение регрессии на отклики, для которых нормальные остатки заведомо неуместны.
Асимптотическая нормальность $\hat\beta$ позволяет использовать критерии Вальда и отношения правдоподобия по аналогии с $t$- и $\chi^2$-тестами гл.~\ref{ch:basics}, хотя конечновыборочные свойства могут отличаться.

Типичные связки:
\begin{center}
\footnotesize
\begin{tabularx}{\linewidth}{@{}>{\raggedright\arraybackslash}X>{\raggedright\arraybackslash}X>{\raggedright\arraybackslash}X>{\raggedright\arraybackslash}X@{}}
\toprule
Отклик & Распределение & $g(\mu)$ & Пример \\
\midrule
Непрерывный & $N(\mu,\sigma^2)$ & $\mu$ & МНК \\
Счётный & $Pois(\lambda)$ & $\log\lambda$ & визиты к врачу \\
Бинарный & $Ber(\pi)$ & $\logit\,\pi$ & перелом, отклик \\
\bottomrule
\end{tabularx}
\end{center}
Сравнение вложенных GLM $G=2(L_{ur}-L_r)\sim\chi^2_{k_1}$ при $H_0$.
При $k_1=1$ критерий Вальда и $G$ не эквивалентны (в отличие от $t$ и $F$ в OLS). При расхождении ориентируйтесь на $G$.

Deviance и остатки GLM. Остаточная deviance $D_{\mathrm{res}}=2(L_{\mathrm{sat}}-L_{\mathrm{fit}})$ аналогична $RSS$. При добавлении регрессоров не возрастает. Псевдо-$R^2_{\mathrm{DEV}}=1-D/D_0$, $D_0$~--- модель только с константой.
Остатки Пирсона $r_i^{(P)}=(y_i-\hat\mu_i)/\sqrt{V(\hat\mu_i)}$, для логистики $V(\mu)=\mu(1-\mu)$.
Deviance-остатки $r_i^{(D)}=\mathrm{sign}(y_i-\hat\mu_i)\sqrt{d_i}$, где $d_i$~--- вклад $i$-го наблюдения в $D_{\mathrm{res}}$.
Стандартизованные остатки deviance используют для выбросов и диагностики. Для бинарных $y$ точечная проверка нормальности остатков малоинформативна~--- смотрите бины по $\hat\mu$ или кривые калибровки.
Deviance играет роль $RSS$ в сравнении вложенных моделей. При $H_0$ статистика $G=2(L_{ur}-L_r)$ асимптотически $\chi^2$~--- тот же принцип, что у частичного $F$ в OLS, но через правдоподобие.

Логистическая регрессия и пробит. Для $y_i\in\{0,1\}$ логит-модель. $\pi_i=P(y_i=1\mid x_i)=\frac{e^{x_i^\top\beta}}{1+e^{x_i^\top\beta}}$.
Пробит $\pi_i=\Phi(x_i^\top\beta)$, где $\Phi$~--- функция стандартного нормального распределения, связь $g(\pi)=\Phi^{-1}(\pi)$.
Логit и пробit дают близкие вероятности в центральной области. Пробit толще в хвостах (экстремальные $P$ меняются быстрее).
Логit удобнее интерпретировать через отношение шансов $\mathrm{OR}=\exp(\beta_j)$. Пробit исторически связан с гауссовой латентной переменной $y^*=x^\top\beta+\varepsilon$, $\varepsilon\sim N(0,1)$.
Выбор связи редко меняет качество подгонки сильно. Важнее спецификация $X$.
Логистическая регрессия~--- естественное продолжение логлинейных моделей для таблиц сопряжённости (гл.~\ref{ch:anova}). Вместо сравнения частот в ячейках мы моделируем $\logit\pi$ как функцию предикторов.

Отношение шансов. При увеличении $x_j$ на~1 шанс умножается на $e^{\beta_j}$.
Проверка $H_0\colon\beta_j=0$. Критерий Вальда, $G$ или score.
Для проверки общей калибровки упоминают критерий Хосмера--Лемешоу. Наблюдения группируют по декилям $\hat\pi$, сравнивают наблюдаемые и ожидаемые частоты $y=1$ (на малых $n$ и при непрерывных $\hat\pi$ критерий спорен. Предпочтительны калибровочные кривые и Brier score).

Качество классификации. ROC-кривая (TPR и FPR при варьировании порога). Площадь под кривой (AUC). На рис.~\ref{fig:reg-roc}~--- типичный вид для бинарного GLM.
AUC измеряет ранжирующую способность модели, но не заменяет проверку калибровки. Модель может хорошо упорядочивать риски и систематически завышать или занижать вероятности.

\begin{figure}[htbp]
    \centering
    \includegraphics[width=\FigWidth]{\FigDir/chapter-5/logit.png}
    \caption{ROC-кривая логистической модели.}
    \label{fig:reg-roc}
\end{figure}

Полная разделимость классов ведёт к бесконечным $\hat\beta$ и огромным стандартным ошибкам. Помогает регуляризация Фирта или штраф.
Порог 0{,}5 не обязан быть оптимален по бизнес-метрике (чувствительность/ специфичность).

Poisson-регрессия и отрицательное биномиальное. Для счётного $y$. $\log\mathbb{E}(y\mid x)=x^\top\beta$, $\mathbb{E}(y\mid x)=e^{x^\top\beta}$.
Poisson предполагает $\mathbb{D}(y\mid x)=\mathbb{E}(y\mid x)$ (экидисперсия).
При $\mathbb{D}(y\mid x)>\mathbb{E}(y\mid x)$ (overdispersion) стандартные ошибки Poisson-GLM занижены.
Отрицательное биномиальное $\mathbb{D}(y\mid x)=\mu+\alpha\mu^2$, $\alpha>0$~--- параметр сверхдисперсии. При $\alpha\to 0$ получаем Poisson.
Оценивают $\hat\alpha,\hat\beta$ МП и тестируют $H_0\colon\alpha=0$ критерием отношения правдоподобия.
Квазиправдоподобие с масштабом $\phi=\sum (y_i-\hat\mu_i)^2/(n-k-1)$ даёт робастные ошибки без полной NB-спецификации.
Underdispersion встречается реже (hurdle/zero-inflated модели).
Для данных \texttt{DVISITS} с избытком нулей типична цепочка. МНК (неадекватен), Poisson-GLM (заниженные ошибки при overdispersion), NB или hurdle-модель.
Содержательный выбор зависит от того, отражают ли нули истинное отсутствие событий или смесь двух подпопуляций.

Сравнение GLM $D_{\mathrm{res}}=2(L_{\mathrm{sat}}-L_{\mathrm{fit}})$. $AIC=-2L+2(k+1)$, $BIC=-2L+\ln n\,(k+1)$, $AICc$~--- поправка при малых $n$.
Разница AIC/BIC $\le 2$ у соседних моделей часто считают несущественной.
Устойчивая дисперсия QML для счётных данных. $\widehat{\cov}(\hat\beta)=(X^\top W X)^{-1}(X^\top D X)(X^\top W X)^{-1}$, $W=\diag(\mu_i)$, $D=\diag((y_i-\mu_i)^2)$.
GLM предполагает правильную спецификацию $g$ и распределения. Для тяжёлых хвостов счётных данных рассматривают zero-inflated модели.
Пропуски MCAR/MAR требуют imputation или моделей с неполными данными. MNAR без модели механизма пропусков~--- систематическое смещение.
Информативные пропуски кодируют индикатором $I[x_j=\mathrm{NA}]$ и заполняют $x_j$ константой вне диапазона.
В прикладных отчётах полезно явно указывать семейство распределения и связь, метод оценивания и сравнения моделей ($G$, AIC), меру качества, релевантную задаче (AUC и калибровка для классификации, средняя ошибка прогноза для счётных данных).

\section{Критерий J, робастные ошибки и расширения GLM}

Формулы и критерии этого раздела дополняют основной текст карточками критериев и деталями оценивания GLM.

Критерий Дэвидсона--Маккиннона для невложенных моделей. Пусть конкурируют спецификации $M_1$ и $M_2$ с прогнозами $\hat y$ и $\hat{\hat y}$ на одной выборке.
Вложите прогноз соперника как дополнительный регрессор:
\[ y = X_1\beta_1 + \gamma_1\hat{\hat y}+\varepsilon,\qquad y = X_2\beta_2 + \gamma_2\hat y+\varepsilon.
\]
Проверьте $H_{01}\colon\gamma_1=0$ и $H_{02}\colon\gamma_2=0$ ($t$- или $F$-критерий).
Таблица решений совпадает с критерием $J$ в основной главе. Карточка для одного вспомогательного теста:

\begin{critcard}{Вспомогательная регрессия Дэвидсона--Маккиннона ($M_1$ против $M_2$)}
Выборка & $(x_i,y_i)$, $i=1,\ldots,n$; $\hat{\hat y}_i$~--- прогнозы $M_2$ \\
Постановка & $y_i = x_{i1}^\top\beta_1 + \gamma_1\hat{\hat y}_i + \varepsilon_i$ \\
$H_0$ & $H_{01}\colon \gamma_1=0$ (дополнительная информация $M_2$ не нужна) \\
Статистика & $t=\hat\gamma_1/\mathrm{SE}(\hat\gamma_1)$ \\
Нулевое расп. & $t_{n-k_1-2}$ (при нормальных ошибках и корректной спецификации) \\
p-value & двустороннее; аналогично для $\gamma_2$ в регрессии с $\hat y$ \\
Интерпретация & оба $H_0$ приняты~--- модели неразличимы; одно отвергнуто~--- вторая предпочтительнее; оба отвергнуты~--- обе неполны
\end{critcard}

Гетероскедастичность. Критерий Бройша--Пагана и робастная ковариация Уайта.

\begin{critcard}{Критерий Бройша--Пагана (гомоскедастичность остатков МНК)}
Выборка & остатки $\hat\varepsilon_i$ регрессии $y$ на $X$ \\
$H_0$ & $\mathbb{D}(\varepsilon_i\mid X)=\sigma^2$ (гомоскедастичность) \\
Статистика & $LM = n\,R^2_{\hat\varepsilon^2}$, где $R^2_{\hat\varepsilon^2}$~--- $R^2$ регрессии $\hat\varepsilon_i^2$ на столбцы $X$ \\
Нулевое распр. & $\chi^2_k$ при $H_0$ ($k$~--- число предикторов без константы) \\
p-value & $P(\chi^2_k \geq LM)$; малый $p$~--- признак гетероскедастичности
\end{critcard}

При гетероскедастичности $\hat\beta$ остаётся несмещённым, но $\widehat{\cov}(\hat\beta\mid X)=\sigma^2(X^\top X)^{-1}$ неверна.
Оценка Уайта (HC0):
\[ \widehat{\cov}_{\mathrm{HC0}}(\hat\beta\mid X) = (X^\top X)^{-1} X^\top \diag(\hat\varepsilon_i^2)\, X (X^\top X)^{-1}.
\]
Модификации HC1--HC4 (MacKinnon, Cribari--Neto) улучшают конечновыборочные свойства. На практике часто используют HC3 при малых $n$.
Диагональ может задаваться $\hat\sigma^2$, $\hat\varepsilon_i^2$, $\hat\varepsilon_i^2/(1-h_{ii})$ и т.\,д.

Расстояние Кука. Расстояние Кука измеряет суммарное влияние $i$-го наблюдения на все подогнанные значения:
\[ D_i=\frac{\sum_{j=1}^n(\hat y_j-\hat y_{j(i)})^2}{(k+1)\hat\sigma^2} =\frac{\hat\varepsilon_i^2}{(k+1)\hat\sigma^2}\cdot\frac{h_{ii}}{(1-h_{ii})^2}, \]
где $\hat y_{j(i)}$~--- предсказание модели, обученной без $i$-го пункта, $h_{ii}$~--- рычаг (диагональ hat matrix $H$).
Пороги на практике. $D_i>1$, $D_i>4/n$, $D_i>3\bar D$, визуальный анализ $D_i$ от $\hat y_i$.
Большой $D_i$ при малом остатке~--- точка тянет наклон. Большой остаток при умеренном $h_{ii}$~--- локальная аномалия $y$.

Оценивание GLM. Scoring Фишера. В GLM $\hat\beta$ находят МП итерационно.
Scoring Фишера на шаге $t$:
\[ \beta^{(t+1)}=\beta^{(t)}+I^{-1}(\beta^{(t)})\,U(\beta^{(t)}), \]
$U(\beta)=\partial\ln L/\partial\beta$~--- score, $I(\beta)=-\mathbb{E}\,\partial^2\ln L/(\partial\beta\partial\beta^\top)$ ~--- информационная матрица Фишера.
Для экспоненциального семейства scoring совпадает с методом Ньютона--Рафсона по $\beta$. На практике используют одну и ту же итерацию с подстановкой $I(\hat\beta)$ после сходимости.
Асимптотически $\hat\beta\sim N(\beta,I^{-1}(\beta))$. Доверительные интервалы и критерий Вальда опираются на $I^{-1}(\hat\beta)$.

Нелинейность связи. LOESS и дробные полиномы. LOESS (локально взвешенная регрессия) для проверки линейности логита по $x_j$. Сглаженная доля успехов $\bar y_{\mathrm{sm}}(x_{ji})$ строится ядерным сглаживанием, затем $\bar\ell_{\mathrm{sm}}(x)=\ln\frac{\bar y_{\mathrm{sm}}}{1-\bar y_{\mathrm{sm}}}$.
График $\bar\ell_{\mathrm{sm}}(x_j)$ должен быть близок к прямой. Систематическая кривизна~--- сигнал добавить $x_j^p$ или заменить связь.

Дробные полиномы (fractional polynomials). Для каждого $x_j$ перебирают степени из набора $S=\{-2,-1,-0.5,0,0.5,1,2,3\}$ (при $p=0$~--- $\ln x_j$), сравнивают модели с одной и двумя степенями по правдоподобию. Если двухчлен не лучше линейной~--- оставляют $x_j$.

Порог классификации в логистической регрессии. Класс $\hat y=1$ при $\hat\pi(x)\geq p_0$.
По умолчанию $p_0=0{,}5$, но для медицинских тестов и контроля качества порог подбирают по целевой чувствительности, специфичности или стоимости ошибок (ROC и матрица ошибок).
Модель может хорошо ранжировать риски (высокий AUC) и плохо калибровать вероятности при $p_0=0{,}5$.

Poisson-GLM. Overdispersion, QML и сэндвич. При $\mathbb{E}(y\mid x)=\mu$ и $\mathbb{D}(y\mid x)=\mu$ (экидисперсия) стандартные ошибки Poisson-GLM состоятельны.
При $\mathbb{D}(y\mid x)=\mu+\alpha\mu^2$, $\alpha>0$ (overdispersion) они занижены. $\hat\beta$ остаётся состоятельным при верной связи $\log\mu=x^\top\beta$.

Диагностика. Отношение $\chi^2/\mathrm{df}_{\mathrm{res}}$ (или $\chi^2$ Пирсона) существенно больше~1. График остатков Пирсона против $\hat\mu$ показывает воронку.

Квазиправдоподобие (QML). При корректном $\mu$,
\[ \widehat{\cov}_{\mathrm{QML}}(\hat\beta) = (X^\top W X)^{-1}(X^\top D X)(X^\top W X)^{-1}, \]
$W=\diag(\mu_i)$, $D=\diag((y_i-\mu_i)^2)$.

Робастная сэндвич-оценка (под произвольную дисперсию $\omega_i$):
\[ \widehat{\cov}_{R}(\hat\beta) = \Bigl(\sum_i \mu_i x_i x_i^\top\Bigr)^{-1} \Bigl(\sum_i (y_i-\mu_i)^2 x_i x_i^\top\Bigr) \Bigl(\sum_i \mu_i x_i x_i^\top\Bigr)^{-1}.
\]
Лечение overdispersion. Отрицательная биномиальная модель, $H_0\colon\alpha=0$ критерием отношения правдоподобия. Quasi-Poisson с масштабом $\phi$ для робастных SE.

Пропуски в данных. MCAR, MAR и информативные.

\begin{itemize}
  \item MCAR (missing completely at random). Вероятность пропуска
        не зависит от данных. Complete-case анализ несмещён, но теряет $n$.
  \item MAR (missing at random). Пропуск в $x_j$ может зависеть
        от других наблюдённых столбцов, но не от скрытого $y$ при фиксированных $x$. Уместны imputation (mice, Amelia), EM, available-case для $X^\top X$.
  \item MNAR. Механизм пропуска зависит от ненаблюдаемого~---
        смещение без модели отбора. Нужны индикаторы и содержательная гипотеза.
\end{itemize}
Если пропуск информативен (отказ отвечать, признак неприменим), создают dummy $I[x_j=\mathrm{NA}]$ и подставляют в $x_j$ константу вне диапазона наблюдений.

\section{Практика и семинар}
\label{sec:practice-ch5}

Практика семинаров 7--8. От регрессионного разминки к GLM для визитов к врачу.

Семинар~7. Разогрев и шумовая регрессия. На семинаре~7 сначала отрабатывают МНК на синтетике. Генерируют $y=x_1+x_2+\varepsilon$ и отдельно вектор \texttt{noise} $\sim N(0,1)$.
Модель только по шуму (\texttt{y \textasciitilde\ noise}) даёт высокий $R^2$ и формально значимый коэффициент при чистом шуме~--- иллюстрация того, что подгонка и p-value на той же выборке не доказывают наличие связи.
Добавление \texttt{noise} в модель с настоящими $x_1,x_2$ не улучшает сравнение по LM-тесту вложенности. Лишний предиктор не нужен.

Далее~--- данные о зарплате и \texttt{looks}. Dummy-кодирование (группы относительно среднего \texttt{looks}), проверка предпосылок, VIF, диагностика остатков и влиятельных точек.
Содержательно это тот же маршрут, что пример красота и оценка курса в основной главе, но с полным циклом отбора признаков.

Эксперимент Фридмана (1983). На лекции~7 при $n\approx k$ пошаговая регрессия на матрице чистого шума находит значимые предикторы, а $t$- и $F$-критерии на той же выборке становятся антиконсервативны.
Семинарный разминка с \texttt{noise}~--- миниатюрная демонстрация той же идеи. Не интерпретируйте p-value без проверки, что сигнал не случаен.

Семинар~8. Число визитов к врачу (DVISITS). Данные Cameron--Trivedi (\texttt{dvi.txt}). Счётный отклик \texttt{DVISITS}, предикторы здоровья и социоэкономики.
Рекомендуемый маршрут анализа:
\begin{enumerate}
  \item Разведка. Гистограмма \texttt{DVISITS}, pairplot с
        \texttt{ILLNESS}, \texttt{ACTDAYS}, \texttt{INCOME}. Исключить редкие уровни (напр., \texttt{FREEPOOR} с малым числом визитов).
  \item МНК и Poisson-GLM на одной формуле. МНК даёт отрицательные
        прогнозы. Poisson~--- первый адекватный семейство.
  \item Диагностика экидисперсии. В \texttt{statsmodels} смотрят
        $\texttt{Pearson chi2 / Df Residuals}$ (или deviance/df).
        Значение $\gg 1$~--- overdispersion, стандартные ошибки Poisson занижены.
  \item Отрицательная биномиальная GLM. $\mathbb{D}(y\mid x)=\mu+\alpha\mu^2$.
        Сравнить AIC/deviance с Poisson. При необходимости упростить формулу (отбор по p-value и критерию $G$).
  \item Сравнение предсказаний. Диаграмму рассеяния \texttt{DVISITS} и $\hat\mu$,
        ранговая корреляция (Kendall). МНК и NB на тех же предикторах.
  \item GLOW500 (опционально). Логистическая регрессия для
        \texttt{FRACTURE}, OR, ROC и выбор порога классификации.
\end{enumerate}
Для счётных данных цепочка МНК $\to$ Poisson $\to$ NB~--- последовательное уточнение семейства, а не выбор одной формулы навсегда.

\section{Контрольные вопросы по теме}

\begin{enumerate}
    \item Запишите МНК-решение $\hat\beta$ и интерпретируйте $R^2$ для связи красоты и оценки курса.
    \item Сформулируйте идею доказательства теоремы Гаусса--Маркова.
    \item Чем частичный $F$-критерий отличается от $t$-критерия для одного $\beta_j$?
    \item Как связаны $h_{ii}$, $\hat\varepsilon_i$ и $D_i$?
    \item Чем отличается доверительный интервал для $\mathbb{E}(y_0\mid x_0)$
          от предсказательного для $y_0$?
    \item Интерпретируйте $\exp(\hat\beta_{\mathrm{SMOKE}})$ в логистической
          модели для GLOW.
    \item Когда Poisson-GLM для \texttt{DVISITS} предпочтительнее OLS?
          Когда нужна отрицательная биномиальная модель?
    \item Чем логит-модель отличается от пробита? Зачем упоминают Хосмера--Лемешоу?
    \item Почему при сравнении GLM рекомендуют критерий отношения правдоподобия,
          а не только Wald?
\end{enumerate}
